% ============================================================
% mpi_body.tex — Multiple PD/PI Leadership Plan
% Required only for MPI submissions. Delete from proposal_package.tex if
% this is a single-PI project.
%
% NIH requirements: describe (1) why MPI is appropriate, (2) each PI's
% distinct role, (3) shared responsibilities, and (4) conflict resolution.
% ============================================================

{\noindent\fontsize{12}{14}\selectfont\bfseries Multiple PD/PI Leadership Plan}

\vspace{3pt}

This project will be conducted using a Multiple Principal Investigator (MPI)
leadership model to integrate complementary expertise in
\placeholder{domain 1 — e.g., clinical science} and
\placeholder{domain 2 — e.g., computational methodology}.
This structure is essential for achieving the proposed aims, which require both
\placeholder{what PI 1 uniquely provides} and \placeholder{what PI 2 uniquely
provides}. Leadership is shared between \placeholder{PI 1 name} and
\placeholder{PI 2 name}, whose roles are clearly delineated, non-overlapping,
and jointly critical to the project's success.

The MPI approach is justified by the scientific scope of the project.
\placeholder{1--2 sentences explaining why no single investigator could
independently provide sufficient depth across the required domains, and how
the MPI model ensures balanced leadership and integrated decision-making.}

\textbf{\placeholder{PI 1 Name, credentials} (\placeholder{role label;
e.g., Contact PI, Lead Content Expert})} will \placeholder{describe PI 1's
primary responsibilities: what they lead, what decisions they own, and what
outputs they are accountable for. Include their authority domain.
3--4 sentences.}

\textbf{\placeholder{PI 2 Name, credentials} (\placeholder{role label;
e.g., Lead Methodologist})} will \placeholder{describe PI 2's primary
responsibilities: what they lead, what decisions they own, and what outputs
they are accountable for. Include their authority domain.
3--4 sentences.}

The two PIs will share responsibility for
\placeholder{list shared responsibilities: e.g., strategic planning, personnel
supervision, joint manuscripts}. Decisions will be made collaboratively by
consensus whenever possible. For domain-specific issues,
\placeholder{describe the decision-making split by domain}.
If consensus cannot be reached, the Contact PI will make the final
determination, with rationale documented.

Communication will be maintained through \placeholder{describe meeting
cadence: e.g., weekly standing PI meetings and ad hoc meetings during
critical phases}. All project materials will be maintained in shared,
version-controlled repositories. In the event of unresolved conflict,
\placeholder{describe institutional mediation procedure at PI's institution}.
